\section*{Appendix}

\subsection*{Simplifying expressions}

To simplify an expression, we can use the following local simplifications:
\begin{itemize}
    \item $x + 0 \mapsto x$
    \item $x - x \mapsto 0$
    \item $-(-x) \mapsto x$
    \item $1 \cdot x \mapsto x$
    \item $0 \cdot x \mapsto 0$
    \item $\frac{x}{x} \mapsto 1$
    \item $\frac{x}{1} \mapsto x$
    \item $\frac{x}{-y} \mapsto -\frac{x}{y}$
    \item $\frac{-x}{y} \mapsto -\frac{x}{y}$
    \item $x^0 \mapsto 1$
    \item $x^1 \mapsto x$
    \item constant propagation
\end{itemize}

\subsection*{DFS for linear elimination order}

In this setting, DFS means that we choose one available linear elimination, apply
it, and continue recursively on the reduced system until no further linear
elimination is possible. Only then do we backtrack and try a different
elimination choice from an earlier step.

\begin{algorithm}[H]
\caption{DFS search for a linear elimination order}
\label{alg:dfs-linear-elimination}
\begin{algorithmic}[1]
\Require Polynomial system $P$
\Ensure A reduced system $Q$ together with a substitution map $\Sigma$
\Function{DFSLinearElimination}{$P, \Sigma$}
    \State $L \gets \{(p, x) \mid p \in P,\ p \text{ is linear in } x\}$
    \If{$L = \emptyset$ or $P = \emptyset$}
        \State \Return $(P, \Sigma)$
    \EndIf
    \State $B \gets (P, \Sigma)$
    \ForAll{$(p, x) \in L$}
        \State solve $p = 0$ for $x$ to obtain $x \gets f(\hat{x})$
        \State $P' \gets P$ with $x \gets f(\hat{x})$ substituted in every equation
        \State remove zero equations from $P'$
        \State $\Sigma' \gets \Sigma \cup \{x \mapsto f(\hat{x})\}$
        \State $R \gets \Call{DFSLinearElimination}{P', \Sigma'}$
        \If{$R = (\emptyset, \Sigma'')$ for some $\Sigma''$}
            \State \Return $R$
        \EndIf
        \If{$R$ is better than $B$}
            \State $B \gets R$
        \EndIf
    \EndFor
    \State \Return $B$
\EndFunction
\end{algorithmic}
\end{algorithm}

Here ``better'' means that the resulting system has fewer
remaining equations, and if the number of remaining equations is the same, then we prefer the system with more free variables.